\section{Projection Operators}

Projection operators are a bit weird. In C, the only "projections" are the numbers 0 and 1. So it doesn't particularly tell us a lot about what it is generalizing.
1 
2 Q and P are equivalent, so this means that projections survive equivalence
3 They will be projections in a more general sense, where you can project non-orthogonally
5
6 
7 Exercises on the sums of projections 
8 
10 The key idea here is that projections kind of encode the same properties as certain subspaces. Certain operations on projections correspond to certain operations on subspaces, not unlike characteristic functions and subsets. Indeed, the analogy is uncanny. Let chiA be the characteristic function of the set A, and pY be the projection on subspace Y:

Intersection of subsets A and B, has char function chiA * chiB
Intersection of subspaces Y and Z, has projection pZ pY 

Union of subsets A and B has char function chiA + chiB - chiA * chiB
Adding two subspace Y and Z has projection pY + pZ - pY pZ 

So in a sense, characteristic functions encode the logic of subsets. And projections encode the "logic" of subspaces, this i the beginning of the field of quantum logic!

\begin{exercise}{1}
fill
\end{exercise}
\begin{proof}
fill
\end{proof} 

\begin{exercise}{2}
fill
\end{exercise}
\begin{proof}
fill
\end{proof} 

\begin{exercise}{3}
fill
\end{exercise}
\begin{proof}
fill
\end{proof} 

\begin{exercise}{5}
fill
\end{exercise}
\begin{proof}
fill
\end{proof} 

\begin{exercise}{6}
fill
\end{exercise}
\begin{proof}
fill
\end{proof} 

\begin{exercise}{7}
fill
\end{exercise}
\begin{proof}
fill
\end{proof} 

\begin{exercise}{8}
fill
\end{exercise}
\begin{proof}
fill
\end{proof} 

\begin{exercise}{10}
fill
\end{exercise}
\begin{proof}
fill
\end{proof} 
